\label{exp:ceb01}

\subsubsection*{Постановка}
\textbf{Физический вопрос.}
Воспроизводят ли целочисленные инварианты кубооктаэдра~\latGCKshort\ ($N_{\mathrm{shell}}(k)=10k^2+2$, $\lfloor\caBitBudget\rfloor=9$, $N_\varphi=13$, $N_{\mathrm{HB}}=63$) структурные параметры молекулярной биологии (шаг стекинга B-ДНК, объём нуклеосомы $147\ \mathrm{bp}$, разбиение Румера $8+(12+1)=21$, спираль Полинга--Кори $3{,}6_{13}$ и замки \texttt{C2H2}/\texttt{Ig}), и укладывается ли полипептидная цепь длины~$L$ ко-трансляционно за линейное число операций $O(L) \le 36 L$ без экспоненциального перебора Левинталя?

\textbf{Что проверяем.}
\begin{enumerate}[label=\textbf{(\alph*)},leftmargin=*,itemsep=0.25em]
\item Точное совпадение числа пар оснований нуклеосомного кора с кумулятивным шаром кубооктаэдра $N_{\mathrm{ball}}(3) = 1+12+42+92 = 147\ \mathrm{bp}$ и числа смысловых классов кода с полусферой $N_{\mathrm{shell}}(2)/2 = 21$.
\item Геометрию правой $\alpha$-спирали $3{,}6_{13}$ ($n_\alpha = 18/5 = 3{,}6$, $h_\alpha = 4\caKappaGeom a_0 \approx 1{,}497\ \text{\normalfont\AA}$, трансмембранный стержень $L_{\mathrm{TM}} = 18 h_\alpha \approx 26{,}94\ \text{\normalfont\AA}$) и складчатого $\beta$-листа ($h_\beta = \frac{28}{3}\caKappaGeom a_0 \approx 3{,}492\ \text{\normalfont\AA}$).
\item Строгую верхнюю границу числа элементарных операций ко-трансляционного фолдера $N_{\mathrm{ops}}(L) \le d\cdot\geomNeighborTwelve\cdot L = 36\,L \in \mathrm{P}$ и распознавание замков \texttt{C2H2} (петля $\geomNeighborTwelve=12$) и \texttt{Ig} (петля $N_{\mathrm{HB}}=63$).
\end{enumerate}

\textbf{Рабочая гипотеза.}
Все кристаллографические эталоны B-ДНК, нуклеосомы (PDB \texttt{1AOI}) и вторичных структур Полинга--Кори воспроизводятся с точностью $<0{,}4\%$, а ко-трансляционный фолдинг на пазах~$K_4$ работает за строго линейное время $O(L)$.

\textbf{Наблюдаемые.}
$h_{\mathrm{stack}}(\text{B-DNA})$, $N_{\mathrm{nuc}}$, $N_{\mathrm{classes}}$, $h_\alpha$, $P_\alpha$, $L_{\mathrm{TM}}$, $h_\beta$ и удельная трудоёмкость фолдинга $N_{\mathrm{ops}}/L$ (табл.~\ref{tab:exp-ce-b-01}).

\textbf{Критерий согласия с теорией.}
Точное совпадение дискретных инвариантов ($147$, $21$, $12$, $63$); отклонение линейных шагов спиралей от рентгеноструктурных эталонов $<0{,}5\%$; выполнение $N_{\mathrm{ops}}/L \le 36$ (тип~B).

\subsubsection*{Теоретическая часть}
В главах~\ref{ch:bio-dna-barcode}--\ref{ch:bio-o1-folding} показано, что азотистые основания и пептидные звенья лежат в плоском гексагональном срезе $\{111\}$ ($120^\circ$), а каждый тетраэдрический узел~$\text{C}_\alpha$ ($K_4$, $\cos\theta=-1/3$) имеет ровно $d=3$ дискретных паза ($\alpha,\beta,\gamma$), выбираемых локально по числу замкнутых водородных мостов~$E_{\mathrm{HB}}$.

\textbf{Проверка размерностей.}
Все пространственные шаги спиралей выражаются через $\caKappaGeom a_0$ и имеют размерность длины~$\mathrm{L}$, а структурные счётчики и алгоритмическая удельная сложность безразмерны:
\[
  [h_{\mathrm{stack}}] = [h_\alpha] = [P_\alpha] = [L_{\mathrm{TM}}] = [h_\beta] = [\caKappaGeom\,a_0] = \mathrm{L}\ (\text{\normalfont\AA}),
\]
\[
  [N_{\mathrm{nuc}}] = [N_{\mathrm{classes}}] = [n_\alpha] = \left[\frac{N_{\mathrm{ops}}}{L}\right] = 1,\qquad
  [\Delta G_{\mathrm{bind}}] = [n_{\mathrm{HB}}\,E_{\mathrm{HB}}] = \mathrm{M\,L^2\,T^{-2}}.
\]

\subsubsection*{Ход эксперимента}
Воспроизведение: \texttt{python verify\_genome\_printer.py} (модули \texttt{dna.py}, \texttt{genetic\_code.py}, \texttt{folding.py}, \texttt{locks.py} в репозитории \texttt{genome-printer}).

\begin{enumerate}[label=\textbf{Шаг \arabic*:},leftmargin=*,itemsep=0.35em]
\item Построить 3D-дуплекс B-ДНК длины $147\ \mathrm{bp}$ и проверить штрихкод большой бороздки.
\item Проверить биективность инволюции Румера $\text{U}\leftrightarrow\text{G}, \text{A}\leftrightarrow\text{C}$ между Октетом~I ($8$) и Октетом~II ($8$).
\item Выполнить ко-трансляционную укладку тестовых полипептидов и измерить $N_{\mathrm{ops}}/L$ и $\Delta G_{\mathrm{bind}}$.
\end{enumerate}

\begin{table}[htbp]
\centering
\footnotesize
\setlength{\tabcolsep}{3.5pt}
\caption{Сверка структурных инвариантов ДНК, хроматина и $O(1)$-фолдинга белка (CE-B-01)}
\label{tab:exp-ce-b-01}
\begin{tabular}{@{}llrrr@{}}
\toprule
Наблюдаемое & Формула из носителя & Модель & Рентген / PDB & Откл. \\
\midrule
Шаг стекинга $h_{\mathrm{stack}}$ B-ДНК, \text{\normalfont\AA} & $\lfloor\caBitBudget\rfloor\caKappaGeom a_0 = 9\caKappaGeom a_0$ & $3{,}368$ & $3{,}38 \pm 0{,}02$ & $-0{,}36\%$ \\
Вместимость нуклеосомы $N_{\mathrm{nuc}}$, bp & $N_{\mathrm{ball}}(d) = 1+12+42+92$ & $147$ & $147$ (\texttt{1AOI}) & $0$ \\
Смысловые классы кода ($20+\text{Stop}$) & $N_{\mathrm{shell}}(2)/2 = 42/2$ & $21$ & $21$ & $0$ \\
Виток $\alpha$-спирали $n_\alpha$, ост./виток & $(\geomNeighborTwelve+|\square|)/d_{\mathrm{out}}(E) = 18/5$ & $3{,}60$ & $3{,}60$ & $0{,}00\%$ \\
Подъём $\alpha$-спирали $h_\alpha$, \text{\normalfont\AA} & $d_{\mathrm{in}}(\square)\caKappaGeom a_0 = 4\caKappaGeom a_0$ & $1{,}497$ & $1{,}500$ & $-0{,}22\%$ \\
Трансмембранный стержень $L_{\mathrm{TM}}$, \text{\normalfont\AA} & $(\geomNeighborTwelve+|\square|)h_\alpha = 72\caKappaGeom a_0$ & $26{,}94$ & $27{,}0 \pm 0{,}3$ & $-0{,}22\%$ \\
Шаг $\beta$-листа $h_\beta$, \text{\normalfont\AA} & $(\lfloor\caBitBudget\rfloor+1/d)\caKappaGeom a_0 = \frac{28}{3}\caKappaGeom a_0$ & $3{,}492$ & $3{,}48 \pm 0{,}02$ & $+0{,}36\%$ \\
Колонка коры ($|\square|=6$ слоёв): клетки / связи & $|\square|d_{\mathrm{soft}} = 42$ / $72+75 = 147$ & $42$ / $147$ & $N_{\mathrm{shell}}(2)$ / $N_{\mathrm{ball}}(3)$ & $0$ \\
Инвариант $\chi_{\mathrm{2D}}$ ($S_1$ / $\mathrm{Tube}_R$ / $\mathrm{Col}_6$) & $V - E + F_\triangle + F_\square^{\mathrm{surf}}$ & $2/0/41$ & $\chi(S^2)/\chi(\mathrm{Cyl})/\Delta\varphi_{\mathrm{disc}}$ & $0$ \\
Трудоёмкость фолдинга $N_{\mathrm{ops}}/L$ & $\le d\cdot\geomNeighborTwelve = 36$ опер./ост. & $18{,}4$ & $\le 36\ (\in\mathrm{P})$ & PASS \\
\bottomrule
\end{tabular}
\end{table}

\subsubsection*{Анализ, расчёт погрешности и выводы}
Дискретные структурные числа ($147$, $21$, $3{,}6$) совпадают с природными эталонами тождественно ($u=0$), а линейные размеры спиралей отличаются от кристаллографических данных Полинга--Кори менее чем на $0{,}36\%$ (тип~B).
Удельное число операций ко-трансляционного фолдера составляет $\approx 18{,}4 < 36$ на аминокислотный остаток.

\textbf{Вывод.}
Молекулярная геометрия ДНК, генетического кода и белкового фолдинга полностью выводится из геометрии кубооктаэдра~\latGCKshort\ и лежит в классе сложности~$\mathrm{P}$.
\textbf{Статус записи:} опыт закрыт по табл.~\ref{tab:exp-ce-b-01}.
